% !TeX root = ../main.tex
% =================================================================================
\chapter{Methodology}\label{chap:methodology}
% =================================================================================

This chapter describes the research methodology of this thesis.
\Cref{sec:methodology:method} introduces the methodology, and \cref{sec:methodology:application} describes how we apply it in this thesis.

% ---------------------------------------------------------------------------------
\section{Research Methodology}\label{sec:methodology:method}
% ---------------------------------------------------------------------------------

Introduce the research methodology you follow and cite its sources.
Many theses that build an artifact (a tool, a method, a language, \ldots) follow Design Science Research, e.g., the three cycle view by \textcite{Hevner2007}:
The \emph{relevance cycle} connects the research to its application domain and provides requirements and field testing.
The \emph{rigor cycle} grounds the research in the existing knowledge base and adds the results to it.
The \emph{design cycle} iteratively builds and evaluates the artifact.
Other methodologies are, e.g., case studies, controlled experiments, surveys, or systematic literature reviews.
Agree on the methodology with your supervisor.

% ---------------------------------------------------------------------------------
\section{Application in This Thesis}\label{sec:methodology:application}
% ---------------------------------------------------------------------------------

Describe the concrete activities you perform and map them to the methodology and to the chapters of your thesis.
For example:
As part of the rigor cycle, we conduct a literature study (\cref{chap:related-work}).
As part of the relevance cycle, we elicit requirements through stakeholder interviews (\cref{chap:requirements}).
As part of the design cycle, we design (\cref{chap:concept}) and implement (\cref{chap:implementation}) our approach in several increments and evaluate it (\cref{chap:evaluation}).
